\chapter*{The Defence of the Old Order: II}
\addcontentsline{toc}{chapter}{The Defence of the Old Order: II}
\setcounter{postnote}{0} % Restarting endnotes numbering for each chapter

\section*{1}
In the politics of Marxism, there is no institution which is nearly as important as the state - so much so that concentration of attention upon it has helped to devalue in Marxist theory other important elements of politics, for instance the cultural elements discussed in the last chapter.

Nevertheless, there is no question that it is right to emphasize the central role of the state in any given social system. Nor of course is the emphasis specific to Marxism: it is also found in other and opposed theories of politics.

Most of these theories, however large their differences from each other may be, have in common a view of the state as charged with the representation of 'society as a whole', as standing above particular and necessarily partial groups, interests, and classes, as having the special function of ensuring both that the competition between these groups, interests, and classes remains orderly and that the 'national interest' should not be impaired in the process. Whether the state discharges these tasks well or not is not here at issue: the point is that most theories of politics do assign the tasks in question to the state, and thereby turn it into the linchpin of the social system.

The starting point of the Marxist theory of politics and the state is its categorical rejection of this view of the state as the trustee, instrument, or agent of 'society as a whole'. This rejection necessarily follows from the Marxist concept of society as class society. In class societies, the concept of 'society as a whole' and of the 'national interest' is clearly a mystification. There may be occasions and matters where the interests of all classes happen to coincide. But for the most part and in essence, these interests are fundamentally and irrevocably at odds, so that the state cannot possibly be their common trustee: the idea that it can is part of the ideological veil which a dominant class draws upon the reality of class rule, so as to legitimate that rule in its own eyes as well as in the eyes of the subordinate classes.

In reality, so the Marxist argument has always been, the state  is an essential means of class domination. It is not a neutral referee arbitrating between competing interests: it is inevitably a deeply engaged partisan. It is not 'above' class struggles but right in them. Its intervention in the affairs of society is crucial, constant and pervasive; and that intervention is closely conditioned by the most fundamental of the state's characteristics, namely that it is a means of class domination - ultimately the most important by far of any such means.

The most famous formulation of the Marxist view of the state occurs in the Communist Manifesto: The executive of the modern state is but a committee for managing the common affairs of the whole bourgeoisie.'\postnote{Revs., p. 69.} But this is not nearly so simple and straightforward a formulation as it has commonly been interpreted to be. In fact, it presents, as do all Marxist formulations on the state, many problems which need careful probing. I do not mean by this that the general perspective is false: on the contrary, I think that it is closer to the political reality of class societies than any other perspective. But it is not a magic formula which renders the interpretation of that reality unproblematic. There is no such formula.

One immediate problem concerns the notion of 'ruling class' in its Marxist usage. In that usage, the 'ruling class' is so designated by virtue of the fact that it owns and controls a predominant part of the means of material and 'mental' production; and that it thereby controls, runs, dictates to, or is predominant in the state as well. But this assumes that class power is automatically translated into state power. In fact, there is no such automatic translation: the question of the relation between class power and state power constitutes a major problem, with many different facets. Even where that relation can be shown to be very close, a number of difficult questions remain to be answered, or at least explored. Not the least of these questions concerns the forms which the state assumes, and why it assumes different forms, and with what consequences.

But the very first thing that is needed is to realize that the relation between the 'ruling class' and the state is a problem, which cannot be assumed away. Indeed, the problem is implicit in the formulation from the Manifesto which I have quoted earlier. For the reference to 'the common affairs of the whole bourgeoisie' clearly implies that the bourgeoisie is a social totality made up of different and therefore potentially or actually conflicting elements, a point which, as was suggested in Chapter II, must be taken as axiomatic for all classes; while 'common affairs' implies the existence of particular ones as well. On this basis, there is an absolutely essential function of mediation and reconciliation to be performed by the state; or rather, it is the state which plays a major role in the performance of that function, the qualification being required because there are other institutions which help in its performance, for instance the parties of the bourgeoisie.

But if the state is to perform this mediating and reconciling function for what are, in effect, different elements or fractions of the bourgeoisie, which have different and conflicting interests, it clearly must have a certain degree of autonomy in relation to the 'ruling class'. In so far as that class is not monolithic, and it never is, it cannot act as a principal to an agent, and 'it' cannot simply use the state as 'its' instrument. In this context, there is no 'it', capable of issuing coherent instructions, least of all in highly complex, fragmented and 'old' societies, where a long process of historical development has brought to predominance a 'ruling class' which harbours many different interests and fractions. This is not to deny the acceptability, with various qualifications, of the term, but only to suggest that the relation of the 'ruling class' to the state is always and in all circumstances bound to be problematic.

The best way to proceed is to ask the simplest possible question, namely why, in Marxist terms, the state should be thought to be the 'instrument' of a 'ruling class'; and an answer to this question should also yield an answer to the question of the validity of the concept of 'ruling class' itself.

Marxists have, in effect, given three distinct answers to the question, none of which have however been adequately theorized.

The first of these has to do with the personnel of the state system, that is to say the fact that the people who are located in the commanding heights of the state, in the executive, administrative, judicial, repressive and legislative branches, have tended to belong to the same class or classes which have dominated the other strategic heights of the society, notably the economic and the cultural ones. Thus in contemporary capitalism, members of the bourgeoisie tend to predominate in the three main sectors of social life, the economic, the political  and the cultural/ideological - the political being understood here as referring mainly to the state apparatus, though the point would apply more widely. Where the people concerned, it is usually added, are not members of the bourgeoisie by social origin, they are later recruited into it by virtue of their education, connections, and way of life.

The assumption which is at work here is that a common social background and origin, education, connections, kinship, and friendship, a similar way of life, result in a cluster of common ideological and political positions and attitudes, common values and perspectives. There is no necessary unanimity of views among the people in question and there may be fairly deep differences between them on this or that issue. But these differences occur within a specific and fairly narrow conservative spectrum. This being the case, it is to be expected, so the argument goes, that those who run the state apparatus should, at the very least, be favourably disposed towards those who own and control the larger part of the means of economic activity, that they should be much better disposed towards them than towards any other interest or class, and that they should seek to serve the interests and purposes of the economically dominant class, the more so since those who run the state power are most likely to be persuaded that to serve these interests and purposes is also, broadly speaking, to serve the 'national interest' or the interests of 'the nation as a whole'.

This is a strong case, easily verifiable by a wealth of evidence. The bourgeois state has tended to be run by people very largely of the same class as the people who commanded the 'private sector' of the economic life of capitalist societies (and for that matter the 'public sector' as well). What I have elsewhere called the state elitehas tended to share the ideological and political presumptions of the economically dominant class. And the state in capitalist society has tended to favour capitalist interests and capitalist enterprise - to put it thus is in fact to understate the bias.

Yet, strong though the case is, it is open to a number of very serious objections. These do not render the consideration of the nature of the state personnel irrelevant. Nor do they in the least affect the notion of the state as a class state. But they do suggest that the correlation which can be established in class terms between the state elite and the economically dominant class is not adequate to settle the issue.

One major objection is that there have been important and frequent exceptions to the general pattern of class correlation. These exceptions have occurred, so to speak, at both the upper and the lower levels of the social scale.

Britain provides the most notable example of the former case. Here, a landowning aristocracy continued for most of the nineteenth century to occupy what may accurately be described as an overwhelmingly preponderant place in the highest reaches of the state apparatus; while Britain was in the same period turning into the 'workshop of the world', the most advanced of all capitalist countries, with a large, solid, and economically powerful capitalist class. The phenomenon was of course familiar to Marx and Engels, who frequently remarked upon it. Thus in an article, The British Constitution', written in 1855, Marx observed that 'although the bourgeoisie, itself only the highest social stratum of the middle classes', had 'gained political recognition as the ruling class', 'this only happened on one condition; namely that the whole business of government in all its details - including even the executive branch of the legislature, that is, the actual making of laws in both Houses of Parliament - remained the guranteed domain of the landed aristocracy';\postnote{SE, p. 282.} and he went on to add that 'now, subjected to certain principles laid down by the bourgeoisie, the aristocracy (which enjoys exclusive power in the Cabinet, in Parliament, in the Civil Service, in the Army and Navy . . .) is being forced at this very moment to sign its own death warrant and to admit before the whole world that it is no longer destined to govern England'.\postnote{Log. cit.}

Of course, the aristocracy was then neither signing its death warrant nor admitting anything of the kind Marx suggested. But even if we leave this aside, there are here very large theoretical as well as empirical questions which remain unanswered concerning the relations of this new 'ruling class' to a landed aristocracy which, according to Marx, remained in complete charge of state power; and also concerning the ways in which the 'compromise' of which Marx also spoke in this article worked themselves out.

The same phenomenon of an aristocratic class exercising power 'on behalf' of capitalism occurred elsewhere, for instance in Germany, and was discussed by Marx and Engels, though in different terms. But it is obvious that this presents a problem for the thesis that the bias of the state is determined by the social class of its leading personnel.

 The problem is at least as great where an important part of the state apparatus is in the hands of members of 'lower' classes. This too has frequently been the case in the history of capitalism. In all capitalist countries, members of the petty bourgeoisie, and increasingly of the working class as well, have made a successful career in the state service, often at the highest levels.\footnote{In the Prison Notebooks, Gramsci asks: 'Does there exist, in a given country, a widespread social stratum in whose economic life and political self-assertion ... the bureaucratic career, either civil or military, is a very important element ? ; and he answered that 'in modern Europe, this stratum can be identified in the medium and small rural bourgeoisie, which is more or less numerous from one country to another', and which he saw as occasionally capable of 'laying down the law' to the ruling class (op. cit., pp. 212, 213).} It may well be argued that many if not most of them have been 'absorbed' into the bourgeoisie precisely by virtue of their success, but the categorization is too wide and subjective to be convincing. In any case, there are telling instances where no such absorption occurred: the most dramatic such instances are those of the Fascist dictators who held close to absolute power for substantial periods of time in Italy and Germany. This will need to be discussed later, but it may be noted here that whatever else may be said about them, it cannot be said about Hitler and Mussolini that they were in any meaningful way absorbed into the German and Italian bourgeoisies and capitalist classes.

In other words, the class bias of the state is not determined, or at least not decisively and conclusively determined, by the social origins of its leading personnel. Nor for that matter has it always been the case that truly bourgeois-led states have necessarily pursued policies of which the capitalist class has approved. On the contrary, it has often been found that such states have been quite seriously at odds with their economically dominant classes. The classic example is that of Roosevelt after his election to the Presidency of the United States\index{America!United States of} in 1932. In short, exclusive reliance on the social character of the state personnel is unhelpful - it creates as many problems as it solves.

The second answer which Marxists have tended to give to the question why the state should be thought to be the 'instrument' of a capitalist 'ruling class' has to do with the economic power which that class is able to wield by virtue of its ownership and control of economic and other resources, and of its strength and influence as a pressure group, in a broad meaning of the term.

There is much strength in this argument too, to which is added by virtue of the growth of economic giants as a characteristic feature of advanced capitalism. These powerful conglomerates are obviously bound to constitute a major reference point for governments; and many of them are multi-national, which means that important international considerations enter into the question. And there is in any case a vital international dimension injected into the process of governmental decisionmaking, and in the power which business is able to exercise in the shaping of that process, because governments have to take very carefully into account the attitudes of other and powerful capitalist governments and of a number of international institutions, agencies and associations, whose primary concern is the defence of the 'free enterprise' system.

Capitalist enterprise is undoubtedly the strongest 'pressure group' in capitalist society; and it is indeed able to command the attention of the state. But this is not the same as saying that the state is the 'instrument' of the capitalist class; and the pressure which business is able to apply upon the state is not in itself sufficient to explain the latter's actions and policies. There are complexities in the decision-making process which the notion of business as pressure group is too rough and unwieldy to explain. There may well be cases where that pressure is decisive. But there are others where it is not. Too great an emphasis upon this aspect of the matter leaves too much out of account.

In particular, it leaves out of account the third answer to the question posed earlier as to the nature of the state, namely a 'structural' dimension, of an objective and impersonal kind. In essence, the argument is simply that the state is the 'instrument' of the 'ruling class' because, given its insertion in the capitalist mode of production, it cannot be anything else. The question does not, on this view, depend on the personnel of the state, or on the pressure which the capitalist class is able to bring upon it: the nature of the state is here determined by the nature and requirements of the mode of production. There are 'structural constraints' which no government, whatever its complexion, wishes, and promises, can ignore or evade. A capitalist economy has its own 'rationality' to which any government and state must sooner or later submit, and usually sooner.

There is a great deal of strength in this 'structural' perspective,  and it must in fact form an integral part of the Marxist view of the state, even though it too has never been adequately theorized. But it also has certain deficiencies which can easily turn into crippling weaknesses.

The strength of the 'structural' explanation is that it helps to understand why governments do act as they do - for instance why governments pledged to far-reaching reforms before reaching office, and indeed elected because they were so pledged, have more often than not failed to carry out more than at best a very small part of their reforming programme. This has often been attributed - not least by Marxists - to the personal failings of leaders, corruption, betrayal, the machinations of civil servants and bankers, or a combination of all these. Such explanations are not necessarily wrong, but they require backing up by the concept (and the fact) of 'structural constraints' which do beset any government working within the context of a particular mode of production.

The weakness of the case is that it makes it very easy to set up arbitrary limits to the possible. There are 'structural constraints' - but how constraining they are is a difficult question; and the temptation is to fall into what I have called a 'hyper- structuralist trap, which deprives 'agents' of any freedom of choice and manoeuvre and turns them into the 'bearers' of objective forces which they are unable to affect. This perspective is bat another form of determinism - which is alien to Marxism and in any case false, which is much more serious.\footnote{See R. Miliband, 'Poulantzas and the Capitalist State' in New Left Review, no. 82, Nov.-Dee. 1973, for a critical assessment of what I take to be an example of this type of determinism, namely N. Poulantzas, Political Power and Social Classes (London, 1973). For a reply, see N. Poulantzas, 'The Capitalist State: A reply to Miliband and Laclau', in New Left Review, no. 95, jan.-Feb. 1976.} Governments can and do press against the 'structural constraints' by which they are beset. Yet, to recognize the existence and the importance of these constraints is also to point to the limits of reform, of which more later, and to make possible a strategy of change which attacks the mode of production that imposes the constraints.

Taken together, as they need to be, these three modes of explanation of the nature of the state - the character of its leading personnel, the pressures exercised by the economically dominant class, and the structural constraints imposed by the mode of production - constitute the Marxist answer to the question why the state should be considered as the 'instrument' of the 'ruling class'.

Yet, there is a powerful reason for rejecting this particular formulation as misleading. This is that, while the state does act, in Marxist terms, on behalf of the 'ruling class', it does not for the most part act at its behest. The state is indeed a class state, the state of the 'ruling class'. But it enjoys a high degree of autonomy and independence in the manner of its operation as a class state, and indeed must have that high degree of autonomy and independence if it is to act as a class state. The notion of the state as an 'instrument' does not fit this fact, and tends to obscure what has come to be seen as a crucial property of the state, namely its relative autonomy from the 'ruling class' and from civil society at large. This notion of the relative autonomy of the state forms an important part of the Marxist theory of the state and was, in one form or another, much discussed by Marx and Engels. The meaning and implication of the concept require further consideration.

\section*{2}
As a preliminary to the discussion of the relative autonomy of the state in Marxist thought, it should be noted that Marx and Engels duly acknowledged, as they could hardly fail to do, the existence through history and in their own times of different forms of state, not only in different modes of production and social formations, but also in the capitalist mode of production itself. In historical terms, they identified forms of state ranging from 'Asiatic' despotism to the Absolutist State\index{Absolutist State}, and including the states of Antiquity and of feudalism; and in their own times, they distinguished between such forms of state as the bourgeois republic, the Bonapartist and Bismarckian states, the English and Czarist forms of state, etc.

On the other hand, these distinctions are never such, for Marx and Engels and for most of their disciples, as to nullify the one absolutely fundamental feature that these states all have in common, namely that they are all class states. This twin fact about the state, namely that it assumes many different and sharply contrasting forms but also that it remains a class state throughout, has been a source of considerable tension in Marxist political thought and debate, and has also had a considerable  influence at one time or another on the strategy of Marxist parties and organizations.

The problem, for revolutionaries, is obvious. The distinctions which clearly must be made between different forms of state - if only because they are too great to be ignored - may lead to a disappearance from view of the central idea of the class nature of all forms of state, whether they are dictatorships or bourgeois democratic regimes. But the attempt to avoid this error of appreciation may also lead - and has indeed led - to the assertion that there is really no difference between different forms of state, and that, to take a specific example, there is really no difference, or at least no really serious difference, between Fascist regimes and bourgeois democratic ones. This is what the Comintern 'line' provided as a perspective during a crucial part of inter-war history, with utterly catastrophic results for working-class movements everywhere, and most notably in Germany. It was only at the Seventh World Congress of the Third International in 1935, long after the Nazis had conquered Germany, that Georgy Dimitrov, speaking on behalf of that disastrous organization, gave the official seal of approval to a drastic change of course, and declared that 'accession to power of fascism is not an ordinary succession of one bourgeois government by another, but a substitution of one state form of class domination of the bourgeoisie - bourgeois democracy - by another form open terrorist dictatorship'; and he also proclaimed that now the toiling masses of the capitalist countries are faced with the necessity of making a definite choice, and making it today, not between proletarian dictatorship and bourgeois democracy, but between bourgeois democracy and fascism'.\postnote{G. Dimitrov, The United Front (London, 1938), pp. 12, 110. Quoted in M. Johnstone, Trotsky and the Popular Front' in Marxism Today, October 1975 p. 311.} In fact, the choice had presented itself years before and had been stated quite clearly inside the ranks of Marxism, most notably by Leon Trotsky.\postnote{See L. Trotsky, The Struggle Against Fascism in Germany (New York 1971)}

Since those days, many different varieties of 'ultra-leftism' have included, as part of their creed, the insistence that there was no real difference between bourgeois democratic regimes and authoritarian ones. That this is an ultra-left deviation seems to me evident. But rather than simply dismiss it as such, one should see it as an expression of the real theoretical and practical problems which the existence of different forms of class state poses to Marxists: not the least of these problems is that, in so far as some of these forms are infinitely preferable to others, choices often have to be made, as Dimitrov said, which involve the defence of bourgeois democratic regimes against their opponents on the right. The terms on which that defence should be conducted has always been a major strategic problem for revolutionary movements, even when they were agreed that the defence itself must be undertaken.

On this latter point, there never was any question for classical Marxism, beginning with Marx and Engels. The latter were not at all taken with bourgeois democracy, and denounced it in utterly uncompromising terms as a form of class domination. Indeed, Marx in The Class Struggles in France wrote of the bourgeois republic of 1848 that it 'could be nothing other than the perfected and most purely developed rule of the whole bourgeois class . . . the synthesis of the Restoration and the July monarchy'.\postnote{SE, p. 89.} Also, a recurrent theme in Marx's writings on the subject is how repressive and brutal this form of state can turn as soon as its upholders and beneficiaries feel themselves to be threatened by the proletariat. With the June days in Paris, the Republic, Marx wrote in the same text, appeared 'in its pure form, as the state whose avowed purpose it is to perpetuate the rule of capital and the slavery of labour'; and 'bourgeois rule, freed from all fetters, was inevitably transformed, all at once, into bourgeois terrorism'.\postnote{Ibid., p. 61.} In the same vein, Marx wrote in The Civil War in France more than two decades later that the treatment meted out to the Communards by the government of Thiers showed what was meant by 'the victory of order, justice and civilization': 'The civilization and justice of bourgeois order comes out in its lurid light whenever the slaves and drudges of that order rise against their masters. Then this civilisation and justice stand forth as undisguised savagery and lawless revenge'.\postnote{FI, p. 226. My italics.}

But when this has been said, it remains true that Marx and Engels saw considerable virtues in bourgeois democratic regimes as compared with other forms of class domination, and notably with Bonapartism, to which Marx in particular devoted much attention. Twenty years of his life, the years of his intellectual maturity, were spent in the shadow, so to speak, of this form of dictatorship; and just as capitalist Britain served as Marx's laboratory for the analysis of the political economy of capitalism, so did France under Louis Bonaparte serve him between 1851 and 1871 as a frame of reference for one kind of capitalist politics.

 The Bonapartist state was distinguished by an extreme inflation and concentration of executive power, personified in one individual, and nullifying the legislative power. This executive power, Marx wrote in The Eighteenth Brumaire of Louis Bonaparte, possessed 'an immense bureaucratic and military organisation, an ingenious and broadly based state machinery, and an army of half a million officials alongside the actual army, which numbers a further half million ; and he further described this state machinery as 'this frightful parasitic body, which surrounds the body of French society like a caul and stops up all its pores\dots\postnote{SE, p. 237.}
 
 There is, in Marx's writings on the Bonapartist state, a strong sense of the life of 'civil society' being stifled and suppressed in ways which the class state of bourgeois democratic regimes cannot in normal circumstances adopt - a sense epitomized in his description of France after thecoup d'etat of December 2: '\dots all classes fall on their knees, equally mute and equally impotent, before the rifle butt.'\postnote{Ibid., p. 236.}
 
 Closely related to this, and crucially important in the approach of Marx and Engels to the bourgeois republican form of state, as compared with other forms of bourgeois state power, was their belief that the former provided the working class with opportunities and means of struggle which it was precisely the purpose of the latter to deny them. One of the most specific and categorical statements connected with this issue occurs in The Class Struggles in France:
 \begin{displayquote}
    The most comprehensive contradiction in the Constitution (of the Second Republic) consists in the fact that it gives political power to the classes whose social slavery it is intended to perpetuate: proletariat, peasants and petty bourgeoisie. And it deprives the bourgeoisie, the class whose old social power it sanctions, of the political guarantees of this power. It imposes on the political rule of the bourgeoisie democratic conditions which constantly help its enemies towards victory and endanger the very basis of bourgeois society. It demands from the one that it should not proceed from political emancipation to social emancipation and from the other that it should not regress from social restoration to political restoration.\postnote{Ibid., p. 71.}
 \end{displayquote}
 
 The main feature of the bourgeois democratic regimes on which Marx and Engels fastened was universal suffrage. Still in The Class Struggles in France, Marx pointed to the problem which this presented, or could present, to the bourgeoisie:
 \begin{displayquote}
    But does the Constitution still have any meaning the moment that the content of this suffrage, this sovereign will, is no longer bourgeois rule? Is it not the duty of the bourgeoisie to regulate the franchise so that it demands what is reasonable, its rule? By repeatedly terminating the existing state power and by creating it anew from itself does not universal suffrage destroy all stability; does it not perpetually call all existing powers into question; does it not destroy authority; does it not threaten to elevate anarchy itself to the level of authority?\postnote{Ibid., p. 127.}
 \end{displayquote}
 
 Marx himself answered that the moment must arrive when universal suffrage was no longer compatible with bourgeois rule, and that this moment had indeed arrived in France with the elections of 10 March 1850. The bourgeoisie was now forced to repudiate universal suffrage, he said, and to confess: 'Our dictatorship has existed hitherto by the will of the people; it must now be consolidated against the will of the people'.\postnote{Ibid., p. 127.} Nevertheless, he also noted, 'universal suffrage had fulfilled its mission, the only function it can have in a revolutionary period. The majority of the people had passed through the school of development it provided. It had to be abolished - by revolution or by reaction'.\postnote{Ibid., p. 134.}
 
 There are many different questions which this conjures up and which will be discussed in Chapter VI. But it may be useful at this point to pursue a little further the issue of universal suffrage and bourgeois democracy as it appeared to Marx and Engels.

Marx himself always held fast to the view that the suffrage, in so far as it helped to sharpen the contradictions of bourgeois society and provided a 'school of development' of the working class, offered definite but limited possibilities to the revolutionary movement.\footnote{In an article on the Chartists published in the New York Daily Tribune of 25 August 1852, Marx went as far as the following: \ . . universal suffrage is the equivalent for political power for the working class of England, where the proletariat forms the large majority of the population, where, in a long, though underground, civil war, it has gained a clear consciousness of its position as a class, and where even the rural districts know no longer any peasants, but only landlords, industrial capitalists (farmers) and hired labourers. The carrying of universal suffrage in England would, therefore, be a far more socialistic measure than anything which has been honoured with the name on the Continent' (The Chartists' in SE, p. 264). Both he and Engels came, particularly after the passage of the Second Reform Act of 1867, to take a rather less sanguine view of the workings of universal suffrage. (See K. Marx and F. Engels, On Britain (Moscow, 1953), passim.)} On the other hand, he also described 'the general suffrage' in class society as 'abused either for the parliamentary sanction of the Holy State Power, or a play in the hands of the ruling classes, only employed by the people to sanction (choose the instruments of) parliamentary class rule once in many years', instead of being adapted, as he thought it would have been by the Commune, 'to its real purpose, to choose by the communes their own functionaries of administration and initiation'.\postnote{FI, p. 251. The quotation is from a first draft of The Civil War in France.}

The question of the suffrage and its uses is also connected with that of the 'transition to socialism'; and Marx, as noted earlier, was willing to allow that there might be some isolated cases where that transition would be achieved by non-violent means, and therefore presumably through electoral means made available by the suffrage. But he was clearly very sceptical about such a process, and took it for granted that it would not be the common pattern. 'The workers' he said in a speech after the last Congress of the First International in 1872, 'will have to seize political power one day in order to construct the new organisation of labour; they will have to overthrow the old politics which bolsterup the old institutions'. But 'we do not claim', he went on, that the road leading to this goal is the same everywhere':
\begin{displayquote}
    We know that heed must be paid to the institutions, customs and traditions of the various countries, and we do not deny that there are countries, such as America\index{America!United States of} and England, and if I was familiar with its institutions, I might include Holland, where the workers may attain their goal by peaceful means. That being the case, we must recognize that in most continental countries the lever of the revolution will have to be force; a resort to force will be necessary one day in order to set up the rule of labour.\postnote{Ibid., p. 324. The Congress was held in The Hague and the speech was delivered in Amsterdam.}
\end{displayquote}

The weight of the argument is unmistakably on the nonpeaceful side of the line. The case of Engels is a little more complicated, the complication arising from pronouncements made in the course of developments, notably the growth of German Social Democracy, which postdate Marx's death. By far the most important such pronouncement is Engels's famous Introduction of 1895 to Marx's Class Struggles in France. That Introduction was subjected to some expurgation before it was published in Vorwarts, the main organ of the German Social Democratic Party; and Engels strongly complained to Kautsky that the expurgated version made him appear as a 'peaceful worshipper of legality quand meme'.\postnote{SW 1950, I, p.109, fn. 1.} But even though Engels had every right to complain that the editing of his text was misleading, the fact remains that the unexpurgated version unquestionably shows a major shift of emphasis from earlier pronouncements of Marx and Engels on the question of the suffrage and its uses. Engels emphatically denied that he was in any way suggesting that 'our foreign comrades' should 'in the least renounce their right to revolution'.\postnote{Ibid., p. 124; p. 665 in SW 1968.} But he went on to say that 'whatever may happen in other countries, the German Social- Democracy occupies a special position and therewith, at least in die immediate future, has a special task'.\postnote{Ibid., p. 124; p. 665 in SW 1968.}

This 'special task' was to maintain and preserve the growth of German Social Democracy and its electoral support:
\begin{displayquote}
    Its growth proceeds as spontaneously, as steadily, as irresistibly, and at the same time as tranquilly as a natural process. All government intervention has proved powerless against it. We can count even today on two and a quarter million voters. If it continues in this fashion, by the end of the century we shall conquer the greater part of the middle strata of society, petty bourgeois and small peasants, and grow into the decisive power in the land, before which all other powers will have to bow, whether they like it or not. To keep this growth going without interruption until it of itself gets beyond the control of the prevailing governmental system, not to fritter away this daily increasing shock force in vanguard skirmishes, but to keep it intact until the decisive day, that is our main task.\postnote{Ibid., p. 125; p. 665 in SW 1968.}
\end{displayquote}

Engels's reference to 'the decisive day' suggests that he had by no means given up the notion that a revolutionary break might yet occur. So does his reference a little later in the text to the likelihood that the powers-that-be would themselves be forced to break through the 'fatal legality' under which Social Democracy thrived, and which it therefore was its bounden duty to maintain. But the shift of emphasis from an earlier perspective is nevertheless quite clear; and whatever Engels himself may have thought, his rejection of notions of'overthrow', even if bound up with specific conjunctural circumstances, was sufficiently decisive to give considerable encouragement to a much more positive view of what was possible under bourgeois democracy than had earlier been envisaged by Marx and Engels. Not, it should be added, that the shift would not have occurred without Engels's text. The text only gave, or appeared to give, legitimation to a process that was inscribed in the evolution of the German labour movement in the conditions of a bourgeois democratic regime even as incomplete and stunted as the German Imperial one.

 In the present context, however, the critical point is that there was no real question for Marxists until the outbreak of the First World War that working-class parties and movements must seize every opportunity of electoral and representative advance offered to them by their respective regimes;\postnote{See Rosa Luxemburg's comment, below, p. 161.} and that they must do whatever was in their power to further the 'democratization' of these regimes. In Marxist thought, bourgeois democracy, for all its class limitations, remained a vastly superior form of state to any existing alternative.

Lenin s changing perspectives, in this connection, are particularly interesting, since they considerably affected later debates. At the time of the Russian Revolution of 1905, he pressed hard the case for proletarian support - indeed for the leadership - of the bourgeois revolution in Russia:
\begin{displayquote}
    In countries like Russia the working class suffers not so much from capitalism as from the insufficient development of capitalism. The working class is, therefore, most certainly interested in the broadest, freest, and most rapid development of capitalism . . . That is why a bourgeois revolution is in the highest degree advantageous to the proletariat. The more complete, determined and consistent the bourgeois revolution, the more assured will the proletariat's struggle be against the bourgeoisie and for socialism.\postnote{Two Tactics of Social-Democracy in the Democratic Revolution, in SWL, p. 76.}
\end{displayquote}

Moreover, Lenin was then very concerned to stress the differences that were to be found within bourgeois democracy itself:
\begin{displayquote}
    There are bourgeois-democratic regimes like the one in Germany, and also like the one in England; like the one in Austria and also like those in America\index{America!United States of} and Switzerland. He would be a fine Marxist indeed, who in a period of democratic revolution failed to see this difference between the degrees of democratism and the difference between its forms, and confined himself to 'clever' remarks to the effect that, after all, this is 'a bourgeois revolution', the fruit of 'bourgeois revolution'.\postnote{Ibid., p. 79.}
\end{displayquote}

Nor was this in any way a purely tactical or debating standpoint. In 1908 he was referring casually to America\index{America!United States of} and Britain as countries 'where complete political liberty exists';\postnote{'Inflammable Material in World Politics', in CWL (1963), vol. 15, p. 186.} and in 1913, writing in the context of the national question, he suggested that 'advanced countries, Switzerland, Belgium, Norway, and others, provide us with an example of how free nations under a really democratic system live together in peace or separate peacefully from each other'.\postnote{'The Working Class and the National Question', ibid., vol. 19, p. 91.}

These glowing remarks should not be taken too literally: Lenin never ceased to make the sharpest distinction between bourgeois democracy and what he envisaged that the proletarian form of democracy could be. But it was only under the impact of the First World War that he adopted a much more undiscriminating stance towards all forms of bourgeois democracy. Thus in the Preface which he wrote in August 1917 to The State and Revolution, he said that
\begin{displayquote}
    the imperialist war has immensely accelerated and intensified the process of transformation of monopoly capitalism into state-monopoly capitalism. The monstrous oppression of the working people by the state, which is merging more and more with the all-powerful capitalist associations is becoming increasingly monstrous. The advanced countries - we mean their hinterland - are becoming military convict prisons for the workers;\postnote{The State and Revolution, in SWL, p. 264.}
\end{displayquote}

and in the pamphlet itself, he specifically said, in connection with his insistence on the need to 'smash' the bourgeois state, that
\begin{displayquote}
    both Britain and America\index{America!United States of}, the biggest and the last representatives - in the whole world - of Anglo-Saxon 'liberty', in the sense that they had no militarist cliques and bureaucracy, have completely sunk into the all-European filthy, bloody morass of bureaucratic-military institutions which subordinate everything to themselves, and suppress everything.\postnote{Ibid., p. 290.}
\end{displayquote}

Even so, he had also noted, a little earlier in the pamphlet, that 'we are in favour of a democratic republic as the best form of state for the proletariat under capitalism', though with the qualification that 'we have no right to forget that wage slavery is the lot of the people even in the most democratic bourgeois republic'.\postnote{Ibid., p. 264.} This was the more or less 'traditional' view. The other, which saw all capitalist countries as 'sunk in the bloody morass of bureaucratic-military institutions' represented a different appreciation, whose 'ultra-leftist' tinges had considerable consequences in terms of policy decisions later. As noted earlier, it was the same trend of thought which led to the adoption, or which at least served as a justification for the adoption, of the 'class against class' Comintern policies of the 'Third Period', in which all bourgeois regimes, of whatever kind, were assimilated for strategic purposes under the same rubric, with results that materially contributed to the Nazi conquest of power in Germany.

There are of course many different reasons why this assimilation was easily accepted by the overwhelming majority of members of Communist parties and organizations. One general reason, as I also already noted earlier, has to do with an attitude of mind to which Marxists have been prone. This is the belief that because A and B are not totally different, they are not really different at all. This error has not only been made in relation to the state, but in other contexts as well, with damaging effects. More specifically, there is a permanent Marxist temptation to devalue the distinction between bourgeois democratic regimes and authoritarian ones. From the view that the former are class regimes of a more or less repressive kind, which is entirely legitimate, it has always been fairly easy for Marxists to move to the inaccurate and dangerous view that what separates them from truly authoritarian regimes is of no great account, or not 'qualitatively' significant. The temptation to blur the distinction has been further enhanced by the fear that to do otherwise would make more difficult an intransigent critique of the class limitations and inherent shortcomings of bourgeois democracy; and by the fear that it would conceal the fact that bourgeois democracy can be turned with the assent and indeed encouragement of ruling classes into authoritarian or Fascist regimes.

Different forms of state have different degrees of autonomy. But all states enjoy some autonomy or independence (the terms are used interchangeably here) from all classes, including the dominant classes.

The relative autonomy of the state was mainly acknowledged by Marx and Engels in connection with forms of state where the executive power dominated all other elements of the state system - for instance the Absolutist State\index{Absolutist State}, or the Bonapartist or Bismarckian one. Where Marx and Engels do acknowledge the relative autonomy of the state, they tend to do so in terms which sometimes exaggerate the extent of that autonomy. Later Marxist political thought, on the contrary, has usually had a strong bias towards the underestimation of the state's relative autonomy.

What this relative autonomy means has already been indicated: it simply consists in the degree of freedom which the state (normally meaning in this context the executive power) has in determining how best to serve what those who hold power conceive to be the 'national interest', and which in fact involves the service of the interests of the ruling class.

Quite clearly, this degree of freedom is in direct relation to the freedom which the executive power and the state in general enjoy vis-a-vis institutions (for instance parliamentary assemblies] and pressure groups which represent or speak for either the dominant class or the subordinate ones. In this sense, the relative autonomy of the state is greatest in regimes where the executive power is least constrained, either by other elements within the state system, or by various forces in civil society. Of such regimes, Bonapartist France was the example with which Marx was most familiar. Its extreme version in capitalist society has been Fascism in Italy and Nazism in Germany, with many other less thorough examples in other capitalist societies. But it may be worth stressing again that all class states do enjoy some degree of autonomy, whatever their form and however 'representative' and 'democratic' they may be. The very notion of the state as an entity separate from civil society implies a certain distance between the two, a relation which implies a disjunction. The disjunction, which may be minimal or substantial, can only come to an end with the disappearance of the state itself, which in turn depends on the disappearance of class divisions and class struggles; and most likely a lot else as well.

I noted earlier that Marx and Engels discussed the state's relative independence mainly in connection with regimes where the executive power was exceptionally strong. This makes it appear that it is only in such regimes that they thought the state was relatively independent. More important, the way they discussed that relative independence in the regimes they did consider is somewhat confusing, in so far as it involves the use, particularly by Engels, of a concept of 'equilibrium' between contending social forces, which 'equilibrium' is supposedly provided by the state.

In The Eighteenth Brumaire, Marx contrasts as follows Bonapartism and the regimes which had preceded it: 'Under the Restoration, Louis Philippe and the parliamentary republic ', what Marx called 'the bureaucracy', and which here stands for the state, 'was the instrument of the ruling class, however much it strove for power in its own right'. 'Only under the second Bonaparte', Marx goes on, 'does the state seem to have attained a completely autonomous position. The state machine has established itself so firmly vis-a-vis civil society that the only leader it needs is the head of the Society of December 10\dots{}' (i.e. Bonaparte).\postnote{SE, p. 238.}

This would appear to suggest the complete independence of the state power from all social forces in civil society. But Marx goes on to say that 'the state power does not hover in mid-air', and that Bonaparte 'represents a class, indeed he represents the most numerous class of French society, the small peasant proprietors'.\postnote{Ibid., p. 238.} What the notion of 'representation' means here is not at all clear, and the discussion of the nature and status of Bonapartism which follows is very weak: on the one hand, Bonaparte is described as 'the executive authority which has attained power in its own right, and as such he feels it to be his mission to safeguard "bourgeois order".\postnote{Ibid., p. 245.} But the strength of this bourgeois order lies in the middle class. He therefore sees himself as the representative of the middle class and he issues decrees in this sense'.\postnote{Ibid., p. 246.} On the other hand, Marx also writes that 'the contradictory task facing the man explains the contradictions of his government, the confused and fumbling attempts to win and then to humiliate first one class and then another, the result being to array them all in uniform opposition to him'.\postnote{Ibid., p. 246.} This obviously understates and indeed obscures the quite specific class role which Bonaparte and his regime were called upon to play.

Twenty years later, in The Civil War in France, the significance of Bonapartism is much more clearly articulated. Marx describes the Second Republic which was overthrown by the coup d'etat of 2 December 1851, as a 'regime of avowed terrorism and deliberate insult toward the "vile multitude'"; and he then goes on to describe how the 'party of Order' paved the way for the Bonapartist dictatorship:
\begin{displayquote}
    The restraints by which their own divisions had under former regimes still checked the state power were removed by their union; and in view of the threatening upheaval of the proletariat, they now used that state power mercilessly and ostentatiously as the national war-engine of capital against labour. In their uninterrupted crusade against the producing masses they were, however, bound not only to invest the executive with continually increased powers of repression, but at the same time to divest their own parliamentary -stronghold - the National Assembly - one by one, of all its own means of defence against the executive. The executive, in the person of Louis Bonaparte, turned them out. The natural offspring of the 'party-of-Order' republic was the Second Empire.\postnote{FI, p. 207.}
\end{displayquote}

Marx described the Empire as 'the only form of government possible at a time when the bourgeoisie had already lost, and the working class had not yet acquired, the faculty of ruling the nation'.\postnote{Ibid., p. 208.} We may leave aside the notion that the French bourgeoisie had, in 1851, lost the faculty of ruling the nation. Also, the formulation lends itself to the view that the state somehow comes 'in between' the contending classes. But the context makes it absolutely clear that Marx meant precisely the opposite.\postnote{Ibid., p. 208.} 'Imperialism', he wrote, here meaning the type of regime represented by the Second Empire, 'is, at the same time, the most prostitute, and the ultimate form of the state power which nascent middle-class society had commenced to elaborate as a means of its own emancipation from feudalism, and which full-grown bourgeois society had finally transformed into a means for the enslavement of labour by capital'.

There is here no doubt about the class nature and the class function of the Napoleonic regime. On the other hand, the notion of the state as serving to maintain (or to bring about) a certain 'equilibrium' between warring classes occurs in the observations of Marx and Engels on the subject of Absolutism,\postnote{See P. Anderson, Lineages of the Absolutist State\index{Absolutist State} (London, 1974).} and was also used by Engels to describe both the Second Empire and the Bismarckian state.

In his The Origin of the Family, Private Property and the State (1884), Engels wrote of the state that: 'it is, as a rule, the state of the most powerful, economically dominant class, which, through the medium of the state, becomes also the politically dominant class, and thus acquires new means of holding down and exploiting the oppressed class.'\postnote{SW 1968, p. 587.}

Having given the state of antiquity, the feudal state and the 'modern representative state' as examples of the class state, he then goes on:
\begin{displayquote}
    By way of exception, however, periods occur in which the warring classes balance each other so nearly that the state power, as ostensible mediator, acquires, for the moment, a certain degree of independence of both. Such was the absolute monarchy of the seventeenth and eighteenth centuries, which held the balance between the nobility and the class of burghers; such was the Bonapartism of the First, and still more of the Second French Empire, which played off the proletariat against the bourgeoisie and the bourgeoisie against the proletariat. The latest performance of this kind ... is the new German Empire of the Bismarck nation: here capitalists and workers are balanced against each other and equally cheated for the benefit of the impoverished Prussian cabbage junkers.\postnote{Ibid., p. 588.}
\end{displayquote}

The point here is not that the characterization is historically  erroneous, but that it departs considerably from the classical theory of the state associated with Marx and Engels. In the formulation above, the state not only acquires a very high degree of independence, albeit 'by way of exception', but this independence appears to free it from its character as a class state: it seems to have become what might be described as a 'state for itself'. In a famous letter written in 1890, Engels strengthens this impression in his reference to 'the interaction of two unequal forces': 'On the one hand, the economic movement, on the other, the new political power, which strives for as much independence as possible, and which, having been once established, is endowed with a movement of its own.'\postnote{F. Engels to C. Schmidt, 27 October 1890, ibid., p. 696.}

In the same passage, Engels explicitly retains the primacy of the 'economic movement': but his formulations do not point to the class identity of the relatively independent state. Engels never did in fact suggest that the state was a 'state for itself', or that it could be. But his theorization of the relative autonomy of the state is, in class terms, nevertheless evidently inadequate and indeed fairly misleading. As was already noted earlier, the relative independence of the state does not reduce its class character: on the contrary, its relative independence makes it possible for the state to play its class role in an appropriately flexible manner. If it really was the simple 'instrument' of the 'ruling class', it would be fatally inhibited in the performance of its role. Its agents absolutely need a measure of freedom in deciding how best to serve the existing social order.

This 'problematic' has the major advantage of helping to explain a crucial attribute of the state in capitalist society, namely its capacity to act as an agency of reform.

Reform has been a major characteristic of capitalist regimes - not surprisingly since reform has been a sine qua non of their perpetuation. What is perhaps less obvious is that it is the state upon which has fallen the prime responsibility for the organization of reform. Power-holders inside the state system have been well aware of the responsibility, and have acted upon that awareness, not because they were opposed to capitalism, but because they wanted to maintain it.

But to act as the organizers of reform, power-holders have needed some elbow room, an area of political manoeuvre in which statecraft in its literal sense could be exercised. What to concede and when to concede - the two being closely related - are matters of some delicacy, which a ruling class, with its eyes fixed on immediate interests and demands, cannot be expected to handle properly. Power-holders themselves may fail but their chances are better and in this respect at least - the defence of the capitalist economic and social order - their record has been fairly successful, though account must here also be taken of the various weaknesses, errors, and difficulties of their opponents.

Still, the point needs to be stressed that much, if not most, of the reform which power-holders have organized in capitalist societies has generally been strongly and even bitterly opposed by one or other fraction of the 'ruling class', or by most of it. Nor is that opposition to reform altogether 'irrational'. No doubt, the resistance to reform on the part of an economically and socially dominant and privileged class is in the long run all but certain to lead to great trouble, instability, revolt, and possibly overthrow. In this sense, such resistance is 'irrational'; and a systematic and comprehensive resistance on principle to all reform is in any case stupid and eccentric. But from the point of view of the class or classes concerned, resistance to reforms organized by power-holders, in so far as that resistance is selective and flexible, cannot be taken as being necessarily 'irrational'. After all, power-holders may well miscalculate and statecraft can go wrong. It is, for instance, possible to argue that there are occasions and circumstances where reform, far from stilling discontent, will encourage demands for more, and further raise expectations: the phenomenon is familiar. Also, even if the general point holds that reform must in the long run be accepted if a social order is to have any chance to perpetuate itself, the price to be paid in the short run is often real and unpalatable. It is nonsense to say, as is often said on some parts of the Marxist left, that reform does not 'really' affect the 'ruling class'. The latter's members squeal much more than is usually warranted. But the squealing is on the other hand rather more than mere sham: the sense of being adversely affected and constrained is real; and this is quite often an accurate reflection of the concrete impact of this or that measure and action of the state.

Ihe fact that the state is the organizer of reform, often against the wishes of large sections of the 'ruling class' is one reason why the bourgeoisie as a whole has not, in 'normal' circumstances, shown any great craving for dictatorship of the Fascist  or authoritarian type. Its members want a state which is strong enough to impose 'law and order', to contain challenge, and to ensure stability. But the authoritarian or Fascist form of state has, from the point of view of an economically and socially privileged class, the great disadvantage not only of placing a vast amount of power in executive hands to the detriment of other elements in the state, but also of making it much more difficult for hitherto powerful bourgeois elements to exercise a restraining and controlling influence on the people now in charge of the state power. In such situations, the exercise of state power can become dangerously unpredictable and much too 'individualized'.

Undoubtedly, there are circumstances when the bourgeoisie, or large sections of it, will choose this option, notwithstanding its disadvantages. In The Eighteenth Brumaire Marx noted that 'once the bourgeoisie saw "tranquillity" endangered by every sign of life in society, how could it want to retain a regime of unrest, its own parliamentary regime, at the head of society?\dots The parliamentary regime leaves everything to the decision of majorities, why then should the great majority outside parliament not want to make the decisions';\postnote{SE, p. 190.} and he went on to say that
\begin{displayquote}
    by branding as 'socialist' what it had previously celebrated as 'liberal', the bourgeoisie confesses that its own interest requires its deliverance from the peril of its own self-government; that to establish peace and quiet in the country its bourgeois parliament must first of all be laid to rest; that its political power must be broken in order to preserve its social power intact; that the individual bourgeois can only continue to exploit the other classes and remain in undisturbed enjoyment of property, family, religion and order on condition that his class is condemned to political insignificance along with the other classes\dots\postnote{Ibid., p. 190. My italics. See also pp. 224-5.}
\end{displayquote}

This is very well said. But it is precisely because authoritarian and Fascist regimes entail at least the danger of 'political insignificance' for people who do have a strong sense of their own significance that they must hesitate to opt for this alternative. In a letter to Marx on 13 April 1866, Engels wrote that 'Bonapartism is after all the real religion of the modern bourgeoisie.'\footnote{K. Marx-F. Engels, Werke (Berlin, 1965), vol. 31, p. 208. 'I see ever more clearly', Engels goes on, 'that the bourgeoisie is not capable of ruling directly, and that where there is no oligarchy, as there is in England, to take on the task of leading the state and society in the interests of the bourgeoisie for a proper remuneration, a Bonapartist semi-dictatorshp is the normal form; it takes in and the big material interests of the bourgeoisie even against the bourgeoisie but leaves it with no part in the process of governing. On the other hand, this ictatorship is itself compelled to adopt against its will the material interests of the bourgeoisie (ibid., p. 208).} But this was not the case when he wrote these words, nor did it become the case later, the authoritarian option is not 'the religion' of the bourgeoisie: it is its last resort (assuming it is an available resort, which is itself a large and interesting question)\postnote{See below, pp. 175-6.} when constitutional rule and a limited form of state appear to be inadequate to meet the challenge from below and to ensure the continuance of the existing system of domination.

\section*{3}
The main points concerning the forms of the capitalist state and its relative autonomy may be illustrated by reference to the functions which it performs. Briefly, four such functions may be istinguished, even though there is much overlap in practice between them: (a) the maintenance of 'law and order' in the territorial area or areas over which the state is formally invested with sovereignty - the repressive function; (b) the fostering of consensus in regard to the existing social order, which also involves the discouragement of 'dissensus' - the ideological- cultural function; (c) the economic function in the broad sense of the term; and (d) the advancement, so far as is possible, of what keld to be the national interest' in relation to external affairs - the international function.

All states perform these functions. But they perform them differently, depending on the kind of society which they serve; and the state in capitalist society also performs them in different ways, depending on a variety of factors and circumstances. The present context, it should be stressed, is that of advanced capitalism: the question of the state in different contexts requires separate treatment.

The repressive function is the most immediately visible one, ui a literal sense, in so far as it is embodied in the policeman, the soldier, the judge, the jailer, and the executioner. But whether immediately visible or not, the state is a major participant in the c ass struggles of capitalist society. In one way or another, it is permanently and pervasively present in the encounter between conflicting groups and classes - these never meet, so to speak, on their own. The state is always involved, even where it is not invoked, if only because it defines the terms on which the encounter occurs by way of legal norms and sanctions.

In the Marxist perspective, and for the reasons which have been advanced in this chapter, the intervention of the state is always and necessarily partisan: as a class state, it always intervenes for the purpose of maintaining the existing system of domination, even where it intervenes to mitigate the harshness of that system of domination.

But the ways in which it intervenes in the affairs of civil society constitute precisely one of the major defining differences between the forms which the capitalist state assumes.

Thus the bourgeois democratic state is so designated because its powers of intervention, among other things, are variously circumscribed and its police powers variously constrained. In the same vein, the authoritarian state has as one of its distinguishing traits the fact that its powers of intervention are far less constrained and its police powers much larger, much less regulated than is the case for the bourgeois democratic state; and the point applies with even greater emphasis to the Fascist-type state proper.

The limits within which the bourgeois democratic state wields its powers of intervention are not rigidly fixed. Police powers, for instance, vary greatly from period to period, and in the same period depending on what and who is involved. The scope and severity of the repressive power of capitalist regimes cannot be overestimated, notwithstanding long traditions of constitutionalism and a hallowed rhetoric of civic freedom. In periods of serious social conflict, this repressive aspect of the bourgeois democratic state is very quickly deployed; and there are large sections of people to whom this aspect of the bourgeois democratic state is very familiar at all times, including times of relative social peace. To these people - the poor, the unemployed, the migrant workers, the non-whites, and large parts of the working class in general - the bourgeois democratic state does not appear in anything like the same guise as it appears to the well-established and the well-to-do.

Even so, there are qualitative differences between such regimes and authoritarian ones. One crucial such difference is that the latter always make it their first task to destroy the defence organizations of the working class - trade unions, parties, co-operatives, associations, and so on. Bourgeois democratic regimes on the other hand have to accept such organizations as part of their political existence, and are also significantly defined in terms of that acceptance. Such regimes curtail, and are forever seeking to curtail further, the rights and prerogatives which the defence organizations of the working class, and notably trade unions, have managed to acquire over the years. Such attempts at curtailment and erosion are a 'normal' part of the class conflict in which the bourgeois democratic state itself is engaged. So are its constant attempts to integrate, absorb, buy off and seduce the leadership of the defence organizations of the working class. But all this is a very different thing from the actual destruction of independent working-class organizations, which is the hallmark of authoritarian and Fascist regimes (but also of Communist ones, for different reasons which also require separate treatment).

In a somewhat different perspective, the repressive function performed by the capitalist state involves the assurance of 'law and order . The reason why law and order' usually appear in Marxist writing on the subject in inverted commas is not, obviously , because the notions of law and of order are here spurious, but that they are shot through in their conception and in their application with the class connotations imposed upon them by the fact that they are the law and order of particular class societies, applied by particular class states. The 'bias' may be much less visible and specific than is expressed in the popular saying that 'there is a law for the rich and there is a law for the poor though the truth that saying contains should not be overlooked. The point, however, is that a general, pervasive, and powerful set of class premisses and practices affects every aspect of law and order.

This is of more general application. Engels noted that the state performs a number of 'common' services and functions. This is indeed the case. But it performs these 'common' functions in ways which are deeply and inherently affected and even shaped by the character and climate' of the class societies in question. Such matters as law and order, health, education, housing, the environment, and 'welfare' in general are, as all else, not only responsive to but determined, or at least powerfully affected, by the 'rationality' of the system.

This problematic is of exceptional and direct importance in relation to the economic function which the state performs in  capitalist society. State intervention in economic life has always been a central, decisive feature in the history of capitalism, so much so that its history cannot begin to be understood without reference to state action, in all capitalist countries and not only those, such as Japan, where the state was most visibly involved in the development of capitalism.

The state and those who were acting on its behalf did not always intervene for the specific purpose of helping capitalism, much less of helping capitalists, for whom power\hyp{}holders, notably those drawn from a different social class, say the landowning aristocracy, have often had much contempt and dislike. The question is not one of purpose or attitude but of 'structural constraints'; or rather that purposes and attitudes, which can make some difference, and in special circumstances a considerable difference, must nevertheless take careful account of the socio\hyp{}economic system which forms the context of the political system and of state action. Moreover, purposes and attitudes, values and aims, judgements and perspectives, are themselves generally shaped or at least greatly affected by that socioeconomic context, so that what appears 'reasonable' by way of state action (or non\hyp{}action) to power-holders will normally be in tune with the 'rationality' and requirements of the socioeconomic system itself: external and antithetical criteria, in so far as they offend that 'rationality', are by definition 'unreasonable'.

These considerations acquire added significance by virtue of the fact that state intervention has become an ever more pronounced feature of the economic life of capitalism in the last hundred years - and even this formulation runs the risk of greatly understating the extreme acceleration of the process in recent decades, as does the notion of 'intervention' which understates the pervasive and permanent presence of the state, in a multitude of forms, in the economic life of advanced capitalism.

The constantly increasing importance which the state must assume under capitalism is well recognized in classical Marxist writing, and was linked to the accurate prognosis of the ways in which capitalist production would develop. In Capital, Marx had shown that centralization and monopoly were part of 'the immanent laws of capitalistic production itself';\postnote{K. Marx, Capital, I, p. 763.} and some ten years after the publication of Capital, Engels had forecast in Anti-Diihring (1878) that, as a culmination of this process of centralization and monopoly, 'the official representative of capitalist society the state - will ultimately have to undertake the direction of production.'\postnote{F. Engels, Anti-Diihring, p. 380.} However far this might be carried, he also intimated, it would not transform the nature of the state:
\begin{displayquote}
    The modern state, no matter what its form, is essentially a capitalist machine, the state of the capitalists, the ideal personification of the total national capital. The more it proceeds to the taking over of the productive forces, the more does it actually become the national capitalist, the more citizens does it exploit. The workers remain wage-workers -  proletarians. The capitalist relation is not done away with. It is rather brought to a head.\postnote{Ibid., p. 382.}
\end{displayquote}

Engels seems to have had in mind a situation not dissimilar to that created by the growth of joint-stock companies, where 'all the social functions of the capitalist are now performed by salaried employees', and where 'the capitalist has no further social function than that of pocketing dividends, tearing off coupons, and gambling on the Stock Exchange, where the different capitalists despoil one another of their capital'.\postnote{Ibid., p. 381.} On the other hand, the continued existence of 'capitalists' would only be compatible with state ownership of the means of production if these capitalists' were assured of dividends to pocket and coupons to tear off. Engels did not pursue the question, save to say that state ownership could provide no solution to the contradictions of capitalist production, and that 'the harmonizing of the modes of production, appropriation, and exchange with the socialized character of the means of production' could only be brought about 'by society openly and directly taking possession of the production forces which have outgrown all control except that of society as a whole'.\postnote{Ibid., p. 382.}

This formulation obviously raises problems. But it is at any rate clear that for Marx and Engels as well as for a later Marxist tradition, the notion of state intervention in the capitalist process of production is an intrinsic part of the analysis of that process of production itself. Indeed, so concerned were Marxist thinkers after Engels to underline the role of the state in twentieth-century capitalism that they came to speak not only of 'monopoly capitalism' but of 'state capitalism' and of 'state monopoly capitalism', the latter form being the 'official' designation of present-day capitalism by all Communist parties.

'State capitalism' was used by Lenin to denote two different  situations: firstly, that already described by Engels, in which the state plays an ever greater role in the productive process of advanced capitalism. Germany, he wrote in 'Left-Wing' Childishness and the Petty-Bourgeois Mentality of 1918, was "'the last word" in modern large-scale capitalist engineering and planned organisation, subordinated to Junker-bourgeois imperialism ,\postnote{SWL, p. 443. This text, which was published inPravda in May 1918, is not to be confused with 'Left-Wing' Communism - An Infantile Disorder, which was published in 1920.} The second situation was an altogether different one, in which 'state capitalism' is used by Lenin in a much more imprecise and blurred sense, namely the post-revolutionary situation in Russia, when he wanted the Soviet government to foster the growth of capitalist enterprise under the strict supervision of the Soviet government and in accordance with 'national accounting and control'.\postnote{Ibid., p. 445.} State capitalism in this case is characterized by the detailed control which a revolutionary workers' state exercises over capitalist production.

Both these usages seem to me arbitrary and misleading. In the first, German', usage, the notion that state intervention in capitalist production amounts to the state 'taking over' capitalism, which is the connotation that 'state capitalism' tends to convey, is clearly mistaken. In this sense, the closest that capitalism has ever been to 'state capitalism' was in Nazi Germany, which answers well Lenin's (and Bukharin's) description of the process of state intervention, regulation, control, and even dictation: but even here, capitalism under the Nazis did not turn into 'state capitalism': it remained what it had always been, namely a system of production mainly carried on by way of private ownership and control of the predominant part of the means of production, distribution, and exchange, with a much greater apparatus of state direction than hitherto imposed upon that system, though it is also worth remembering that much of that apparatus of direction was itself manned and controlled by people who were part of the traditional German capitalist class.

Nor is there much to be said for the second meaning given to state capitalism in Lenin's usage. Even if the assumption is made that a predominantly capitalist economy can co-exist with and be directed by a revolutionary regime - and it is a very large and to my mind an unwarranted assumption - there would still be no ground for describing such an economic system as 'state capitalism': at most, it could be described as 'state-directed capitalism', or some such. In fact, it would be a capitalist economic system which a revolutionary government would seek to develop, direct, and plan, and into which it might seek to inject some socialist' modes of behaviour.

'State capitalism' has also been used in some Marxist quarters, and by the Chinese Communists, to designate collectivist regimes of the Soviet type, and notably the Soviet Union, mainly as a term of abuse intended to suggest the gap between achievement and promise. But however great the gap may be, and however justified the view that these regimes are not 'socialist' there is no theoretical or practical justification for designating as slate capitalist' regimes in which private ownership and control of the whole or of the largest and most important part of the means of economic activity has been abolished - save as a term of abuse, in which case there is no problem.

As for siate monopoly capitalism', the designation is misleading inasmuch as it tends to suggest something like a symbiotic relation between the contemporary capitalist state and monopoly capitalism. This is inaccurate and opens the way to an over\hyp{}simplified and reductionist view of the capitalist state as the instrument' of the monopolies. In reality, there is the capitalist state on the one hand and 'monopoly capitalism' on the other. The relation between them is close and getting ever closer, but there is nothing to be gained, and much by way of insight to be lost, by a reductionist over-simplification of that relation. 'State monopoly capitalism' does not leave enough space, so to speak, in the relation between the two sides to make the concept a useful one: in so reducing this space, it renders the relation too unproblematic, at any level.

Nevertheiess, the state does of course 'intervene' massively in e of advanced capitalism, and sustains it in a multitude of dlfferent ways which cannot all by any means be labelled 'economic'.\postnote{For instance the cultural realm.} It mainly does so in accordance with the 'rationality' of the capitalist mode of production, and within the constraints imposed upon it by that mode of production. But this is not a simple process either, least of all in a bourgeois democratic and constitutional setting. For in such a setting, the political system a lows pressures to be generated and expressed against and in the state and may turn the state itself into an arena of conflict, with different parts of the state system at odds with each other and thereby reducing greatly the coherence which the state requires to fulfil its functions, 'economic' or otherwise. From this point of view, the authoritarian alternative is also intended  to restore coherence to a state system which has come to reflect rather than subdue the class conflicts and the socio-economic contradictions of civil society. In the conditions of late twentieth-century capitalism, the state in bourgeois democratic regimes is under constant pressure to meet the expectations and demands of the subordinate classes. What is wanted from it is that it should provide and manage - as it alone can - a vast range of collective and public services whose level largely defines the conditions of life for the overwhelming majority of the populations of advanced capitalist countries, who depend upon these services. But against the expectations and demands emanating 'from below' must be set the requirements of capitalist enterprise; and whatever the state does by way of provision and management of services and economic intervention has to run the gauntlet of the economic imperatives dictated by the requirements of the system - and what emerges as a result is always very battered. Hence the contradictory pulls within the state itself, and the need, somehow, to restore its coherence and its capacity to fulfil what is expected of it.

The authoritarian alternative I have referred to is in its extreme forms the ultimate resort of dominant classes, or rather it may be such a resort: the assumption has generally been much too easily made in Marxist writing that dominant classes could simply 'choose' whether to adopt the authoritarian alternative or not. The matter is a lot more complicated than this. But in any case, the bourgeois democratic state is perfectly capable of deploying a vast apparatus of repression, and of resorting to considerable police powers within the constitutional framework in existence. What is involved is the further restriction of civic rights and of the opportunities of organized protest and challenge, and the attempts by the bourgeois state to achieve these restrictions, particularly in times of crisis, is a 'normal', unsurprising part of the civic history of capitalism.

On the other hand, success in this field partly depends on the degree of consent and legitimacy which the state and the social order enjoy, particularly among the subordinate classes themselves; and this links up with a third function of the state that was mentioned earlier, namely its ideological-cultural or persuasive function.

I have already noted that this persuasive function is not primarily carried out by the state in bourgeois democratic regimes, where much or most of it is left to a variety of agencies and institutions which are part of civil society and which largely express the class power of the dominant classes. But it was also noted there that the state is also deeply involved, in such regimes and directly or indirectly, with the business of legitimating the existing social order and with the discouragement of dissidence. This intervention assumes many different forms, which I do not propose to discuss further here, though it may be worth repeating the point that was made in Chapter III that the bourgeois state is ever more closely involved in this form of intervention in the life of civil society; and that the availability of immeasurably more efficient means of communication provides a further spur to such efforts.

Even so, the difference in this realm as in many others between the bourgeois democratic state and the various forms which authoritarianism has assumed in capitalist societies is very great. One of these differences is of course that in the latter case dissident ideas and dangerous thoughts, not only in the realm of politics but in many other realms as well, are forcibly suppressed. Some of this also occurs in bourgeois democratic regimes: but to nowhere near the same extent, nor in a systematic way. A second difference is that the authoritarian state in capitalist regimes itself assumes the main responsibility for the spread of officially approved ideas. It may do so either directly, say by the use of radio, television, and government-sponsored newspapers and other publications; or more indirectly, through rascist-type parties and other organizations; or both. But it is in any case thoroughly involved in the propagation, effectively or not o. ideas which are deemed acceptable and 'functional'- and m the suppression of ideas which are not. The same is also true tor Communist states, which have thoroughly monopolized lorms of communication, through the Party and the various institutions it controls.

The fourth of the functions of the state to which I referred earlier, its international function, raises many different questions with which Marxism has in one way and another always been greatly concerned. This is discussed in the next section.

\section*{4}
From the beginning, Marxism has treated the state as part of a world of states whose relations and actions are deeply influenced and even determined by the fact of capitalist development. In the international as in the national sphere, the socioeconomic context is crucial, particularly because of the supranational character of the capitalist mode of production.

In a famous passage of the Communist Manifesto, Marx and Engels noted the centralizing tendencies of capitalism, as a result of which 'independent, or but loosely connected provinces, with separate interests, laws, governments, and systems of taxation, became lumped together into one nation, with one government, one code of laws, one national interest, one frontier and one customs tariff'.\postnote{Revs., p. 72.} But just as it created the conditions for the centralized, concentrated nation-state, so did capitalism also set in motion powerful supra-national tendencies:
\begin{displayquote}
    The bourgeoisie has through its exploitation of the world market given a cosmopolitan character to production and consumption in every country ... it has drawn from under the feet of industry the national ground on which it stood ... in place of the old local and national seclusion and self-sufficiency, we have intercourse in every direction, universal interdependence of nations . . . The bourgeoisie, by the rapid improvement of all instruments of production, by the immensely faci- lited means of communication, draws all, even the most barbarian, nations into civilization.\postnote{Ibid., p. 71.}
\end{displayquote}

The reference to 'even the most barbarian nations' has a distinctly 'Victorian' ring, but no matter: Marx and Engels were, here as so often elsewhere, accurately projecting the lines of future capitalist development. On the other hand, the processes to which they were pointing have produced intense contradictions, which have multiplied rather than diminished with the passage of time, and have presented major and unresolved problems for Marxist theory and Communist practice.

The most basic of the contradictions in the world-wide development of capitalism is that its expansion did indeed generate 'intercourse in every direction, universal interdependence of nations', and thus created 'one world' in an immeasurably more meaningful sense than had ever been the case before; but it also and simultaneously strengthened, and in many cases engendered, the will to sovereign statehood. On the one hand, there was the economic knitting together of one world; on the other, a heightened drive to political fragmentation by way of the sovereign state in a world of states, each seeking to maximize its power of independent action.

A paramount explanation of this contradiction is that statehood is the absolutely essential condition, though not of course a sufficient one, for the achievement of the aims which those who are able to determine or influence state action may have. It may not be possible, for whatever reason, to achieve these aims with statehood: but it is likely to be even more difficult to achieve them without it.

More specifically, capitalist development and expansion has meant in effect the development of particular national capitalisms, with their respective national states seeking to advance capitalist interests. Obviously, this does not mean that the people in charge of state power necessarily acted with the conscious purpose of serving these interests. What they were doing, in their own view, was to serve the 'national interest', fulfilling their country's 'manifest destiny', spreading civilization and Christianity, serving Queen or Emperor, or whatever. But none of these aims, they also thought, was incompatible with the advancement of national business interests. On the contrary, the advancement of these interests, and their defence against other national interests protected by their own states, was generally felt by power-holders to be congruent and even synonymous with whatever other purposes they had in mind.

The world which Marx and Engels knew in their later years was dominated by large states involved in various forms of rivalry with each other. Some of them, like France, Britain, Russia, and Austria, had had statehood for a long time, and only needed to concern themselves with the preservation of their national and imperial domains, or with further additions to it by way of imperialist expansion. Others, like Germany and Italy, had achieved statehood much more recently and were concerned both with consolidation and with the carving of a share in imperial expansion. One major state, the United States\index{America!United States of}, had had to fight a civil war to retain unitary statehood; and another, the Ottoman Empire, barely held together disparate nationalities.

These were all 'established' states. But the logic which made them cling jealously to wTiat they had, also impelled various national groupings inside one or other of these empires to seek their own statehood. In the nineteenth century, this mainly involved various subject nationalities everywhere. The motivation behind this drive is usually described as 'nationalism'. It is a  convenient term but not an altogether adequate one. For it suggests a specific ideological commitment common to an enormously varied scatter of groups and classes (including varied groups and classes forming part of the same subject nationality), which had in fact very different economic, social, political, cultural aims. 'Nationalism' does not properly cover these different aims, and is too easily turned into a catch-all formula, much too diffuse and imprecise. The term cannot be dispensed with, of course, for it describes some very real drives. But these drives and many others are given much greater precision if they are seen to point, as they all do, to statehood: however diverse the groups and their aims, and whether based on 'nationality' or not, what they do have in common is the will to acquire their 'own' state, to be 'masters in their own house', to achieve independent statehood - often on the basis of a combination of nationalities, ethnic, linguistic, and cultural groupings, depending on particular historical and local circumstances.

So powerful has this drive to statehood been that it has resulted in the proliferation of states to the point where almost 150 of them (a majority of recent creation) are now recognized as 'sovereign' political units, each with the right to a seat in the United Nations. Nor is the power of that drive to statehood at all spent: on the contrary, it is constantly manifesting itself in new places, and has come to pose fairly serious problems to old\hyp{}established states like Britain and France, where demands are heard much more loudly than for a very long time past not merely for greater autonomy but for actual independence of one or other constituent element of the existing 'national' state. There is nothing mysterious about this drive to statehood: it simply marks the recognition that 'sovereignty', however limited it may be, makes possible the fulfilment of aims which may otherwise be unattainable. Whether the aims are good or bad, reasonable or not, is not here in question.

The appeal which statehood has and the support it generates are not confined to any particular class or social group. 'Nationalism' in the nineteenth century and the drive to statehood had exceptional appeal to various elements of the ' national bourgeoisie' of various subject 'nations'. But it was not confined to such bourgeois elements. Nor certainly has it so been confined in the twentieth century. 'Nationalism', more or less strongly seasoned with various other ideological ingredients is available to, and where suitable seized by, any class or group aspiring to independent statehood.

With certain qualifications, classical Marxism was on the whole favourably inclined towards the achievement of independent statehood by subject peoples, though the whole question has become ever more difficult for Marxists to handle in the course of the twentieth century.

The main qualification, in regard to Marx and Engels, is that they quite rightly saw much if not most of the 'nationalism' expressed in their own times as oriented to bourgeois aims indeed as intended, at least in part, to provide an alternative to revolutionary socialism. In this sense, the adoption by the proletariat of a 'nationalism' free from socialist and revolutionary perspectives, and designed to create supra-class allegiances, was naturally seen by them as constituting an instance of 'false consciousness'.

Marx and Engels said in the Communist Manifesto that national differences, and antagonisms between peoples, are daily more and more vanishing, owing to the development of the bourgeoisie, to freedom of commerce, to the world market, to unformity in the mode of production and in the conditions of life corresponding thereto'.\postnote{Ibid., p. 85.} But however much they might hope and strive for the development of proletarian internationalism and solidarity, they continued to think of the nation-state as the basic unit of political life; and to support the right of nations, such as the Irish, the Poles, and the Italians, to national independence and statehood.

However, Marx and Engels also tended to favour larger rather than smaller national units; and it is significant, in this context, that Marx, in The Civil War in France, should have defended the Paris Commune against the accusation that it sought to destroy the unity of France. 'The Communal constitution', he wrote, 'has been mistaken for an attempt to break up into a federation of small states, as dreamed of by Montesquieu and the Girondins, that unity of great nations which, if originally brought about by political force, has now become a powerful coefficient of social production'.\postnote{FI, p. 211.}

This attitude of qualified but quite strong support for the right of 'peoples' to self-determination and statehood is bound to present problems in practice, and has often done so: there are always conflicting considerations which intrude into the question.

Lenin reaffirmed again and again the right of subject peoples  to self-determination and independent statehood: 'In the same way', he wrote, 'as mankind can arrive at the abolition of classes only through a transition period of the dictatorship of the oppressed class, it can arrive at the inevitable integration of nations only through a transition period of the complete emancipation of all oppressed nations, i.e. their freedom to secede.'\postnote{V. I. Lenin, The Socialist Revolution and the Right of Nations to Self- Determination, in SWL, p. 163.} But this did not prevent him from freely admitting that the matter could not be treated in absolute terms. Thus, he noted in the same text that Marx was occasionally taxed with having objected to the 'national movement of certain peoples' (i.e. the Czechs in 1848), and that this was taken to refute 'the necessity of recognizing the self-determination of nations from the Marxist standpoint'. But this, Lenin claimed, was incorrect, 'for in 1848 there were historical and political grounds for drawing a distinction between "reactionary" and revolutionary- democratic nations. Marx was right to condemn the former and defend the latter. The right of self-determination is one of the demands of democracy which must naturally be subordinated to its general interests.'\postnote{Ibid., p. 163, fn. 1.}

Irrespective of the merits of the particular argument, this too is a substantial qualification. But it remains true that Lenin and the Bolsheviks were strongly committed to the right of subject peoples to self-determination. The one major Marxist figure to reject this position was Rosa Luxemburg, who consistently denounced the Tight of self-determination' as hollow phraseology and as a diversionary, corrupting, and self-defeating slogan.

It is important not to drawr false contrapositions here. The point is not that Lenin and the Bolsheviks 'believed' in national independence and statehood for subject peoples, while Rosa Luxemburg did not: to view the controversy in this light is to misunderstand its basis. The point is that Lenin and the Bolsheviks believed that Marxists could not deny the right of subject peoples to independence, and that Marxists who belonged to an oppressor nation, as they did, could do so least of all. This is not the same as 'believing' that independence was always to be encouraged. As for Rosa Luxemburg, she did not of course believe that subject peoples should be denied the right to independence: but that this could only be achieved on the basis of an international socialist struggle that must not be 'diverted' by the acceptance of such slogans as the 'right to self-determination'.

It is a little ironic in the light of experience that Rosa Luxemburg should have denounced so scathingly the Bolsheviks' 'doc- trinaire obstinacy' on this issue, and their adherence to what she called 'this nationalistic slogan'.\postnote{R. Luxemburg, The Russian Revolution, in Rosa Luxemburg Speaks (New York, 1970), p. 380.} In the course of 1917, she warned that acceptance of the right of secession, far from encouraging the revolutionary forces in the subject nations of the Russian Empire, 'supplied the bourgeoisie in all border states with the finest, the most desirable pretext, the very banner of the counterrevolutionary efforts ... By this nationalistic demand they brought on the disintegration of Russia itself, pressed into the enemy's hand the knife which it was to thrust into the heart of the Russian Revolution.'\postnote{Ibid., p. 382.}

In fact, the Bolsheviks themselves became aware of the force of the argument almost as soon as they came to power, and very quickly retreated from their commitment. As early as January 1918, Stalin was telling the Third All-Russia Congress of Soviets that the right of self-determination should be taken to mean 'the right to self-determination not of the bourgeoisie but of the labouring masses of the given nation. The principle of self- determination should be a means in the struggle for socialism and should be subordinated to the principles of socialism.'\postnote{M. Liebman, Leninism under Lenin (London, 1975), p. 273.} With varying degrees of reluctance, or of enthusiasm, this was the view which the Bolsheviks came to adopt, which meant in effect that the right to secession, while formally retained, was for all practical purposes abandoned for the peoples who formed part of Soviet Russia.

On the other hand, the 'right to self-determination' for subject peoples remained a central part of Marxist and Communist thinking in relation to the national question; and it has always figured very large in the anti-colonial and anti-imperialist struggles of this century. However, the qualifications remain and no clear line has ever been established in either Marxist theory or practice to decide where the right to independent statehood is or is not 'justified'. Communist attitudes and policies were long shaped by the internal and external requirements of the Soviet Union; and the contradictory policies adopted by the U.S.S.R. and China over various independence movements for instance Biafra and Bangla Desh - indicate well enough how much practical considerations impinge upon the determination of attitudes to self-determination.

In recent years, as already noted, the question has rapidly crept up on the agenda of 'old' states as well as forming a key Problem for newer' ones. On the whole, Marxist attitudes to demands lor statehood on the part of ethnic, linguistic, and  national forces such as the Scots, the Basques, the Catalans, the Corsicans, the Bretons, the Welsh, etc., have been negative to the point of opposition and hostility. Demands for regional autonomy or for a federal system present few problems to Marxists, who are in favour of it. 'National' claims, on the other hand, do present major problems, of the kind to which Rosa Luxemburg pointed. But as against this, the Leninist stress remains there may come a point, even in 'old states, where demands for statehood on the part of a constituent element of an existing state - e.g. Scotland vis-a-vis the United Kingdom cannot be resisted on principled socialist grounds, or cannot be resisted without resort to suppression, which comes to the same thing. From this point of view, 'nationalism' has proved a much more enduring and therefore a much more difficult problem to confront than early Marxists thought likely; and its emergence as a major problem in 'old' states is a token of the strength of the drive to statehood which was noted earlier. In any case, that drive does not detract but rather reinforces the state as the key unit of political life. States may change and fragment, but the state, with its claims to territorial 'sovereignty' and 'independence' remains.

This fact is not, as might at first sight appear, seriously affected by economic developments in the life of capitalism which are subsumed under the rubric of the multi-national corporation. No doubt, the growth of capitalist enterprise, and the coming into being of the multi-national giants, as was noted earlier, is a major fact in the life of the state, in so far as those who hold power in it have to take the existence of these giants into careful account, to put it very modestly, when they decide upon policy. But there is nothing essentially novel about this; nor is there about the fact that/oreign enterprises, located in any particular national state, can have an important role in the determination of its policies.

This in effect is what the multi-national corporation amounts to in regard to the capitalist state. 'Multi-national' is in any case misleading for most of the enterprises concerned: the largest number of them are United States corporations, with branches in many countries. It is this which makes them 'multi-national', not multi-nationality of ownership or control. Behind the corporations stand their governments: and it is the power which they are able to wield and the iniluence which they are able to exercise which constitute the additional and sometimes the decisive factor of constraint upon the policies and actions of the 'sovereign' nation-state.

But this does not affect the nature of that state. Whatever else in the development of capitalism may require a revision of the Marxist theory of the state, the multi-national corporation does not. On the contrary, the latter merely reinforces the argument that the capitalist state is subject to formidable capitalist constraints - in this instance with an additional international dimension.

At the same time, the point needs to be made that the multinational corporation creates a situation in which the nationstate is seen by some elements at least of the indigenous capitalist class as a necessary instrument of defence against what amounts in effect to foreign interests; while other such elements may seek alliance with these interests rather than defence against them. But in any case, and whetherthis is so ornot, the nature of the state is not thereby brought into question.

\section*{5}
What has so far been said about the Marxist view of the state mainly refers to its nature and role in advanced capitalist societies. The question which this raises is how far if at all Marxist concepts on this subject are appropriate for the analysis of the state in different types of society, namely 'Third World' and Communist ones. Clearly, it is not possible simply to transpose the categories of analysis used for the state in advanced capitalist societies to these different economic, social, and political structures: but what else to do and in what manner to proceed either for 'Third World' societies or for Communist ones is a large problem which, as I noted in the Introduction, has so far not been adequately tackled in Marxist political theory, particularly in relation to Communist societies, where serious Marxist work on the state and politics was for many decades virtually impossible, and where it even now remains peculiarly difficult. The present section purports to do no more than offer certain suggestions about the directions which the required theorization might take.

The first and most obvious feature of the state in both 'Third  World' and Communist societies is a very pronounced inflation of state and executive power, particularly marked in Communist regimes, and with which is usually associated a very high degree of autonomy, at least from the civil society over which the state holds sway. What has been an occasional phenomenon in advanced capitalist societies by way of the extreme inflation of executive power is a common one in these other types of society, though the reasons for this are not the same, or not necessarily the same, in 'Third World' and Communist countries.

In the case of 'Third World' societies, this inflation occurs because social groups which would have an interest in limiting and controlling the power of the state do not have the power or the will to do so; while dominant classes and groups, where they do exist, find it to their advantage to have a strong and repressive state to act on their behalf.

'Third World', as I noted when I first used the term in the Introduction, covers a multitude of specific circumstances and milieux, but in the present context, the countries which are subsumed under this label may be divided into two major categories: those where economically dominant classes, well- entrenched and developed, do exist; and those where they do not. The first category would include the Latin-American\index{America!Latin} continent (with nowadays the obvious exception of Cuba) and the countries of South Asia. The second category would include most countries of Africa\index{Africa}, with the major exception of South Africa, which is an advanced capitalist country. In both categories, foreign capitalist interests constitute an important and in some cases a decisive political element as well as a crucially important economic one.

In the case of countries of the 'Third World' where an economically dominant class exists, or where there is more than one such class, the Marxist analysis of the state presented earlier only requires various adaptations and changes of emphasis rather than fundamental revision. No doubt, there are many aspects of their economic and social structure, their history and political culture, which greatly affect the character and operation of the state in each of them. Differences in economic and social structure produce considerable differences in the nature of class antagonisms and their expression, and in the manner in which the state responds to them. So does the fact of 'under- development' and dependence', and often of recent colonialism, affect the state's weight and role in economic, social and political life. But all these are specificities which, however important and determinant, are susceptible to an analysis that remains identifiably 'Marxist'; in other words to an analysis that is squarely and adequately derived from a Marxist problematic.

The same cannot be said with quite the same confidence about countries that belong to the second category, namely those where an economically dominant class or group, or a number of such classes or groups, did not exist before the establishment of a 'new' state in place of a colonial regime. Of course, it is perfectly possible to point in all such countries to a scatter of local entrepreneurs and traders, alongside the large foreign interests that may exist and which form pockets of large-scale enterprise in an otherwise 'under-developed' context. But the scatter of local entrepreneurs and traders cannot seriously be said to constitute an 'economically dominant class'. Nor can the foreign interests that are present be so designated. This is not to say that these interests are not important politically; or that local business interests, however small the scale of their enterprises, do not represent a point of reference for the state. The point is rather that the element which is absolutely basic in the classical Marxist view of the state, namely an economically dominant class, is not to be found here, in any meaning that makes real economic, social, and political sense. This being so, the question at issue is what the state power in these societies actually 'represents', and what its nature and role may be said to be.

The answer is that, in such societies, the state must be taken mainly to 'represent' itself, in the sense that those people who occupy the leading positions in the state system will use their power, inter alia, to advance their own economic interests, and the economic interests of their families, friends, and followers, or clients. A process of enrichment occurs, which assumes a great number of forms and leads to a proliferation of diverse economic ventures and activities. In this process, a genuine local bourgeoisie may come into being and grow strong, with continuing close connections to the state and its leading members, who are themselves part of that new bourgeoisie.

In such cases the relation between economic and political power has been inverted: it is not economic power which results in the wielding of political power and influence and which shapes political decision making. It is rather political power (which also means here administrative and military power) which creates the possibilities of enrichment and which provides the basis for the formation of an economically powerful class, which may in due course become an economically dominant one. The state is here the source of economic power as well as an instrument of it: state power is a major 'means of production'.

It is an instrument of economic power, not in the sense that those who hold state power serve the interests of an economically dominant class separate from these power-holders and located in society at large; but that those who hold state power use it for their own economic purposes and the economic purposes of whoever they choose. This use of state power assumes many different forms, including of course the suppression of any challenge to the supremacy of what turns in effect into an economically and politically dominant class.

No more than its counterparts anywhere else is this dominant class homogeneous and united. On the contrary, it tends to be very divided and fragmented, with power highly personalized and therefore subject to frequent and violent changes. These societies are poorly articulated in economic and social terms, and therefore in political terms as well. They are in fact de- politicized, with the state itself, often under military rule, assuming a monopoly of political activity through parties and other groupings which are seldom more than bureaucratic shells with very little living substance.

In such circumstances, the state assumes a very high degree of autonomy indeed, and does, almost, become a 'state for itself', or at least for those who command it. There are some qualifications to this, but few. One of these, which may be very considerable, is the existence of foreign interests, which need to be accommodated by the power-holders. On the other hand, the process of accommodation of these foreign interests is not, in many cases, quite as one-sided as it used to be. There are now many 'Third World' countries and states which have a greater freedom of manoeuvre vis-a-vis such interests than used to be the case: there are now fewer 'banana republics' than there were. But if so, it means that the degree of autonomy of the state is commensu- rately increased, and the area of freedom which is available to power-holders similarly enlarged. What they have to fear is not the restraint of their power, but its sudden and usually violent termination by way of a coup engineered by rival groups of people, themselves representing little except their own ambitions and interests.

The question of the state in Communist countries presents very different problems from those which arise in the 'Third World'. There are many different reasons for this, but the most important one is the fact that the economic context is fundamentally different and affects every aspect of the state's role.

In the 'Third World', what the state does is decisively affected by the fact that capitalist enterprise either exists already in one form or other and on a substantial scale, or that it exists on a small scale but can be developed further. This growth of capitalist enterprise is indeed what 'development' largely means in these countries. The 'rationality' of state action in the 'Third World' is determined by this possibility of economic development by way of state encouragement of capitalist enterprise. This state action may not be particularly effective and development may be sluggish or deformed, but that is not here relevant. For all practical purposes, and in terms of what it does, the state has, so to speak, nowhere else to go, no other 'rationality' to follow. The one qualification is that the state may come to be completely dominated by one man and his family and followers, who proceed to stifle and suppress all activity except for such activity as may be of benefit to them. Haiti under Duvalier is a case in point, but such instances are rare and of no great consequence, save for the unfortunate people who live in these countries.

The decisive fact about Communist countries is that they are collectivist regimes in which capitalist enterprise is for all practical purposes non-existent. What there is of it is kept to minimal levels and is positively prevented from growing. It is possible, in these countries, to accumulate a fair amount of money by saving on large salaries, and also to own some property, for instance a house and even a 'second residence'. But it is not possible to make private wealth grow by way of capitalist enterprise.

Just as the possibility of such enterprise decisively shapes state action in the 'Third World', so does its absence shape equally decisively the character and role of the state in collectivist societies. For it means that those who control state power in these societies are subject to 'structural constraints' of a most The Defence of the Old Order: II ill formidable kind, in so far as they cannot direct that power to private capitalist purposes. They may seek and achieve private enrichment, though the scale on which this occurs is fairly limited. And they may well act to the advantage of some classes or groups rather than to that of others - some categories of workers rather than others, or workers as against peasants, or managers as against either, and so on. But this is a very far cry indeed from thepossibilitiesofferedbytheexistenceofacapitalist context, and imposes an altogether different 'rationality' upon those who control collectivist societies, and who are in this sense controlled by the collectivism over which they preside.

It was precisely the fear that the U.S.S.R. under Stalin was moving towards the restoration of capitalism, with the dynamic implications which this would have had, which led the Trotskyist Opposition to warn, from the twenties onwards, of the dangers of a Russian 'Thermidor': and had a capitalist sector been restored in the Soviet Union, by way of a return of a major part of the public sector to private enterprise, the chances are indeed that a Thermidor', in one form or other, would have occurred, and fundamentally reshaped the 'rationality' of the Soviet state.

But in fact, no such restoration occurred. On the contrary, every part of private economic activity was ruthlessly stamped out, most notably by forced collectivization of the countryside. The experience of other Communist countries has in many respects been different from that of the U.S.S.R. But as I have already noted, they are all predominantly and solidly collectivist in their mode of economic organization; and this leaves open the question of the nature and role of the state in this kind of system.

Left critics of Soviet-type regimes have pointed to the very considerable inequalities of power and reward which are to be found in them, and which are sanctioned, maintained, defended, and fostered by an exceedingly powerful state; and they have consequently argued that this state was the instrument ol a new class', 'bureaucratic stratum', 'state bourgeoisie', whose principal purpose, like that of any other 'ruling class', was to maintain and enlarge its power and privileges.

There has over the years been much controversy among Marxist critics of these regimes over whether their rulers did constitute a class or not. The point is obviously of some importance, in so far as the answer to it may provide an initial clue to the degree of cohesion, solidarity, community of purpose, and social basis of these rulers.

However, and as should have been expected, no conclusive answer to the question has ever been returned, or can be. There are Marxists who have said that, because 'the bureaucracy' could not own capitalist property and therefore pass it on to their descendants, they did not form a class; and there are other Marxists who have argued that capitalist property was not the only criterion to be used, and that the privileges which accrued to the people concerned, and from which their descendants could derive advantages of a substantial kind, did mean that 'the bureaucracy' constituted a class.

The word matters less than the substance which it designates; and it is scarcely a matter of serious doubt that those who occupy leading positions in Soviet-type systems do enjoy advantages which are denied to the mass of the population. These advantages may be greater in some of these systems than in others; and greater efforts are made in some cases than in others to reduce such advantages and to de-institutionalize them, at least to some degree. But they clearly continue to form part of Communist life, though they are not of course peculiar to it, and are in fact significantly lower, in material terms, than in other systems. 'Bureaucrats' in collectivist systems are better off than those over whom they rule, and the higher the position a 'bureaucrat' occupies, the better off he is likely to be: but with the probable exception of those at the very top, the pickings would seem to be comparatively modest. Office is an avenue to material wellbeing; but not to great wealth.

This is by no means accidental. It is primarily the result of the absence of opportunities for capitalist enterprise. No doubt there are other factors which account for this aspect of political life in Communist countries, but here is the primary constraint.

This, however, is no refutation of the thesis that the state in these countries is the instrument of a 'power elite', who may be unable to derive vast material advantages from their position in the state system, yet who are able to use that position to appropriate and enjoy tar from negligible such advantages; and who are also able to enjoy power. On this view, the notion of the state as having as its function the advancement of the power and privileges of those who control it, and the repression of those  who challenge it, remains fairly well intact. The 'state bourgeoisie or whatever it is called is not here viewed as a greatly plutocratic class or stratum: but it is nevertheless a privileged class or stratum, and privileged above all in its enjoyment of power.

There is a great deal more in this view than the propagandists of these regimes and other apologists would ever concede. Indeed, the view of Communist regimes as dominated by a 'power elite' located at the top and upper echelons of a pyramid of power seems entirely apposite and sensible. But it is nevertheless entirely inadequate to explain the nature, function, and dynamic of the state in these societies. The reason for this is that it places far too great an emphasis on the purposes and motivations of those who hold power; and it concomitantly grossly underestimates the massive fact constituted by the collectivist context in which they operate and wield their power.

In all Communist countries, the greatest possible emphasis has been placed throughout on economic growth and development, notably of the industrial sector, which was in some cases minimal when the Communists took over. But it was also taken for granted from the beginning that this economic development must occur under the auspices of the state, one of whose absolutely prime functions it was to plan and organize this enterprise and to take all necessary measures for the purpose, including in many cases the ruthless coercion of unwilling populations. No doubt, 'the state' must here be taken to include the Party, whose role in the process has been crucial. But in so far as the Party and state leaders have more or less been the same people, with the Party organization acting as the arm of the Party-state, this presents no particular problem of analysis in the present context.

The point is not that this vast process, amounting to one of the greatest upheavals in human history, and certainly the greatest organized and engineered upheaval of its kind, is incompatible with the 'power and privileges' thesis. There is nothing intrinsically absurd in the idea that many of those who have held power in these regimes were mainly or even wholly moved by the wish to use the state for the aggrandizement of their own power and privileges, whatever the rationalizations they used and believed: the purposes by which all people in power anywhere are moved must be reckoned to be endlessly varied and to span the whole gamut of human motivations, from the noblest to the basest. The point is rather that, whatever the particular motives of the people in power, these motives could only be fulfilled by serving larger purposes as well.

But this is where collectivist constraints play a determinant role in sh aping the functions of the state. For in the absence of an economically dominant capitalist class, and of opportunities for capitalist enrichment by those holding state power, there only remained certain options open - in effect state-fostered economic growth in the broad sense, and including the state- fostered provision of social services and cultural developments; national defence, also of course under the responsibility of the state; and the maintenance of Taw and order', including the massive enlargement of police powers for the purpose of repressing various forms of dissidence.

The same question that was asked earlier in regard to 'Third World' countries arises here: whom and what does this state 'represent' in the discharge of these functions? The best answer would seem to be that it 'represents' no single class or group and is the instrument of no such class or group: the collectivist character of the society precludes it from being such an instrument, for the reasons stated earlier. Instead, the state may be taken to 'represent' the collectivist society or system itself, and to have as its function the service of its needs as these are perceived and defined by those who control the state.

This answer differs fundamentally from that returned by the theorists of the power and privileges' thesis, according to whom the state in Communist societies must be taken to 'represent' the interests of 'the bureaucracy', the 'state bourgeoisie', the 'new class', etc.

It also differs - but not nearly as much - from the designation which was given to the Soviet state in the Khrushchev era, and which has remained the orthodox way of describing it, namely as a 'state of the whole people'. This is unacceptable because it begs too many questions: it has an intensely ideological and apologetic ring, and it assumes, among other things, a quality of representation of the state, in the sense of actually 'representing the will and wishes of the people. But this is an assumption which the nature of the political system makes it impossible to be seriously tested.

'The state' effectively means here the leaders of the Communist Party. For it is the leadership of the Party, and location at  its topmost levels, which confers the power to rule and to use the power of the state for the purpose of ruling. Unlike most countries of the 'Third World' (Mexico being probably the most notable exception), Communist regimes are all distinguished by highly structured and organized mass parties, whose role in the articulation of power in these regimes is absolutely crucial. They are for the most part authentically mass parties, with a membership running, in the case of the U.S.S.R. and China, into many millions. But these mass parties are also distinguished by their pyramidal structure, with an extreme concentration of power at the top of the pyramid.

Some party leaderships may be more autocratic than others, and more personalized. But all these parties are instruments of concentrated power; and the 'democratic centralism' which they all profess has always been a figurative term for an extreme concentration of power and the narrow subordination of the lower organs to the top one. Moreover, Communist parties in these regimes have always, in practice, held a monopoly of political power, whether they were, as in the U.S.S.R., the only legal party, or whether, as in some other Communist countries, other parties were allowed to exist. The latter's existence never seriously impinged upon the 'leading role' of the Communist Party, meaning in effect its complete predominance.

On this basis, those who control the Party also control the state, which is their executive and coercive arm. Nor is the Party-state constrained by social forces - let alone political ones - which are external to it. Of course, party and state leaders must take some account of existing social classes and groupings in determining the policies which they wish to pursue: if they do not, or if they miscalculate, there may be trouble. But their ability to make autonomous decisions, without reference to anyone outside their own restricted circle, is very great indeed and has often been all but unlimited, the extreme example being that of Stalin over a period of a quarter of a century.

This autonomy is constrained, as has already been noted, by the structural context in which Communist societies operate, and this is in many ways very severe. But this constraint does not reduce the power-holders' autonomy vis-a-vis the societies over which they preside, which is the question at issue.

Similarly, Party-state leaders are constrained, save in very unusual cases of total predominance such as Stalin came to enjoy, by internal rivalries and divisions, and the existence at the top of the pyramid of factions and tendencies which may 'represent' particular interests, e.g. the army and defence apparatus, the managers, or managers of some sectors against those of others, and so on; and there may also be divisions based on ideology, differences of generations and experiences, and all the other factors which cause men of power to divide. These are very genuine constraints upon leaders, even upon leaders who have established their ascendency, but who must yet tread warily, and conciliate or disarm opposition, or crush it with the aid of allies, which itself may require concessions to these allies.

But nothing of this fundamentally affects the fact that the state in these systems does have a very high degree of autonomy from society: nowhere is that autonomy substantially constrained by institutions and agencies, either inside the state system or outside, which are in any real sense independent of executive power. 'Freedom', Marx said in the Critique of the Gotha Programme of 1875, 'consists in converting the state from an organ superimposed on society into one thoroughly subordinate to it; and even today state forms are more or less free depending on the degree to which they restrict the "freedom of the state".'\postnote{FI, p. 354.} There is a vast programme of political construction wdiich these formulations propose; and it is a programme to which Communist regimes have not so far seriously addressed themselves.

\printpostnotes